\documentclass[11pt,a4paper]{article}

% =========================================================================
%  RAPPORT HARMELY — Application mobile musicale
%  Theme : Dev mobile UX — palette violet / rose
% =========================================================================

\usepackage[utf8]{inputenc}
\usepackage[T1]{fontenc}
\usepackage[french]{babel}
\usepackage{lmodern}
\usepackage[a4paper,margin=2.2cm,top=2.6cm,bottom=2.4cm]{geometry}
\usepackage{xcolor}
\usepackage{fontawesome5}
\usepackage{tikz}
\usetikzlibrary{shapes.geometric,positioning,calc,arrows.meta}
\usepackage{titlesec}
\usepackage[most]{tcolorbox}
\usepackage{listings}
\usepackage{fancyhdr}
\usepackage{booktabs}
\usepackage{tabularx}
\usepackage{array}
\usepackage{colortbl}
\usepackage{enumitem}
\usepackage{microtype}
\usepackage[bookmarks=true,colorlinks=true,linkcolor=accent,urlcolor=accent,citecolor=accent]{hyperref}

% --- Palette mobile dev UX -----------------------------------------------
\definecolor{primary}{HTML}{1F1147}
\definecolor{accent}{HTML}{7C3AED}
\definecolor{accent2}{HTML}{EC4899}
\definecolor{accentdark}{HTML}{4C1D95}
\definecolor{light}{HTML}{F5F3FF}
\definecolor{midgray}{HTML}{6B7280}
\definecolor{codebg}{HTML}{1A1130}
\definecolor{codefg}{HTML}{E9E6F4}
\definecolor{codekw}{HTML}{D8B4FE}
\definecolor{codestr}{HTML}{FBCFE8}
\definecolor{codecmt}{HTML}{7C7390}

\titleformat{\section}{\sffamily\Large\bfseries\color{primary}}{\thesection}{0.8em}{}[\vspace{-0.4ex}{\color{accent}\rule{\linewidth}{1.2pt}}]
\titleformat{\subsection}{\sffamily\large\bfseries\color{accent}}{\thesubsection}{0.6em}{}
\titlespacing*{\section}{0pt}{1.6em}{0.8em}

\pagestyle{fancy}
\fancyhf{}
\renewcommand{\headrulewidth}{0pt}
\fancyhead[L]{\footnotesize\sffamily\color{midgray}\faIcon{mobile-alt}~Harmely — App musicale}
\fancyhead[R]{\footnotesize\sffamily\color{midgray}Honorine \& Grondin}
\fancyfoot[C]{\footnotesize\sffamily\color{midgray}\thepage}

\lstdefinestyle{jsstyle}{
  backgroundcolor=\color{codebg},
  basicstyle=\ttfamily\footnotesize\color{codefg},
  keywordstyle=\color{codekw}\bfseries,
  stringstyle=\color{codestr},
  commentstyle=\color{codecmt}\itshape,
  showstringspaces=false, breaklines=true,
  frame=none, framesep=8pt, framerule=0pt,
  xleftmargin=12pt, xrightmargin=12pt,
  columns=fullflexible, inputencoding=utf8, extendedchars=true,
  literate={é}{{\'e}}1 {è}{{\`e}}1 {ê}{{\^e}}1 {à}{{\`a}}1 {ô}{{\^o}}1 {î}{{\^\i}}1 {ç}{{\c{c}}}1 {ù}{{\`u}}1 {â}{{\^a}}1 {É}{{\'E}}1 {È}{{\`E}}1 {Ê}{{\^E}}1 {À}{{\`A}}1 {Ô}{{\^O}}1 {Ç}{{\c{C}}}1 {Ù}{{\`U}}1 {Â}{{\^A}}1
}
\lstset{style=jsstyle,language=Java}

\newtcblisting{codebox}[1][]{
  listing only,
  enhanced jigsaw, breakable,
  colback=codebg, colframe=accent,
  arc=2pt, boxrule=0pt, leftrule=3pt,
  left=8pt, right=8pt, top=4pt, bottom=4pt,
  listing options={style=jsstyle,language=Java}, #1
}

\newtcolorbox{infobox}[1][]{
  enhanced, colback=light, colframe=accent,
  arc=2pt, boxrule=0pt, leftrule=3pt,
  left=10pt, right=10pt, top=6pt, bottom=6pt,
  fontupper=\small, #1
}
\newtcolorbox{stepbox}[1][]{
  enhanced, colback=accent!8, colframe=accent2,
  arc=2pt, boxrule=0pt, leftrule=3pt,
  left=10pt, right=10pt, top=6pt, bottom=6pt,
  fontupper=\small, #1
}

\setlist[itemize]{leftmargin=*,itemsep=2pt,topsep=4pt}
\setlist[enumerate]{leftmargin=*,itemsep=2pt,topsep=4pt}

\begin{document}
\pagenumbering{gobble}

\begin{tikzpicture}[remember picture, overlay]
  \fill[primary] (current page.north west) rectangle ([yshift=-9cm]current page.north east);
  \fill[accent] ([yshift=-9cm]current page.north west) rectangle ([yshift=-9.25cm]current page.north east);
  % Logo : smartphone avec notes de musique
  \node[anchor=center] at ([yshift=-4.6cm]current page.north) {%
    \begin{tikzpicture}[scale=1.0]
      % phone outline
      \fill[white,rounded corners=8pt] (-1.1,-1.7) rectangle (1.1,1.7);
      \fill[primary,rounded corners=6pt] (-0.95,-1.3) rectangle (0.95,1.5);
      % speaker
      \fill[white] (-0.18,1.32) rectangle (0.18,1.4);
      % home button
      \fill[white] (0,-1.5) circle (0.1);
      % album art / gradient screen
      \fill[accent2,rounded corners=4pt] (-0.85,0.3) rectangle (0.85,1.3);
      \node[scale=1.4,white] at (0,0.8) {\faMusic};
      % progress bar
      \fill[accent!40] (-0.85,0.1) rectangle (0.85,0.18);
      \fill[accent] (-0.85,0.1) rectangle (0.2,0.18);
      % controls
      \node[scale=0.6,white] at (-0.5,-0.3) {\faStepBackward};
      \node[scale=0.9,white] at (0,-0.3) {\faPlay};
      \node[scale=0.6,white] at (0.5,-0.3) {\faStepForward};
      % swipe right/left hint
      \node[scale=0.9,accent!70] at (-0.55,-0.85) {\faHeart};
      \node[scale=0.9,white!60!accent2] at (0.55,-0.85) {\faTimes};
      % floating notes
      \node[scale=1.0,accent2!80] at (-1.8,1.2) {\faMusic};
      \node[scale=0.8,accent2!60] at (1.7,0.6) {\faMusic};
      \node[scale=0.7,accent!60] at (-1.6,-0.4) {\faMusic};
    \end{tikzpicture}};
  \node[anchor=north, white, font=\sffamily\Huge\bfseries] at ([yshift=-6.7cm]current page.north) {Harmely};
  \node[anchor=north, accent2, font=\sffamily\Large] at ([yshift=-7.65cm]current page.north) {App de découverte musicale};
  \node[anchor=north, white!75!primary, font=\sffamily\small] at ([yshift=-8.4cm]current page.north) {React Native \textbullet{} Firebase \textbullet{} API Spotify};

  \node[anchor=south west, primary, font=\sffamily\small] at ([xshift=2.2cm,yshift=3cm]current page.south west) {\textbf{\textcolor{accent}{Auteurs}}};
  \node[anchor=south west, primary, font=\sffamily\small] at ([xshift=2.2cm,yshift=2.55cm]current page.south west) {HONORINE Kylian};
  \node[anchor=south west, primary, font=\sffamily\small] at ([xshift=2.2cm,yshift=2.25cm]current page.south west) {GRONDIN Angélique};
  \node[anchor=south east, primary, font=\sffamily\small] at ([xshift=-2.2cm,yshift=3cm]current page.south east) {\textbf{\textcolor{accent}{Cadre}}};
  \node[anchor=south east, primary, font=\sffamily\small] at ([xshift=-2.2cm,yshift=2.5cm]current page.south east) {SAE — RT2 TP1};
  \fill[accent2] ([yshift=1.3cm]current page.south west) rectangle ([yshift=1.4cm]current page.south east);
  \node[anchor=south, midgray, font=\sffamily\footnotesize] at ([yshift=0.6cm]current page.south) {IUT Réseaux \& Télécommunications — Université de La Réunion};
\end{tikzpicture}
\newpage

\pagenumbering{arabic}
\setcounter{page}{1}

{\sffamily
\hypersetup{linkcolor=primary}
\renewcommand{\contentsname}{\textcolor{accent}{\faList~} Sommaire}
\tableofcontents
}
\vspace{0.4cm}

\section{Présentation}

\begin{infobox}[title={\faRocket~Pitch produit}]
\textbf{Harmely} est une application mobile de découverte musicale. Inspirée du modèle \textbf{Tinder pour la musique}, elle propose à l'utilisateur des extraits Spotify via un système de swipe \textemdash{} droite pour aimer, gauche pour passer \textemdash{} et conserve l'historique des préférences. Un chat en temps réel permet de partager ses découvertes.
\end{infobox}

\subsection{Stack technique}

\renewcommand{\arraystretch}{1.3}
\begin{tabularx}{\linewidth}{@{}>{\bfseries}p{3.5cm} X@{}}
\toprule
\rowcolor{light} Couche & Technologie \\ \midrule
Front mobile      & React Native (via Expo) \\
Backend           & Firebase Authentication + Firestore \\
Persistance       & AsyncStorage (local) + Firestore (cloud) \\
API tierce        & Spotify Web API (extraits 30\,s) \\
Navigation        & React Navigation (stack) \\
Temps réel        & Firestore \texttt{onSnapshot} \\
\bottomrule
\end{tabularx}

\section{Mise en place de l'environnement}

\begin{enumerate}
  \item Installation de \textbf{Node.js} et \texttt{npm} (gestion des paquets).
  \item Initialisation du projet : \texttt{npm init} génère un \texttt{package.json}.
  \item Installation d'Expo + CLI :
\end{enumerate}

\begin{codebox}[listing options={style=jsstyle,language=bash}]
npm install -g expo-cli
expo init harmely
cd harmely
npm install firebase react-native expo-av expo-font expo-splash-screen
\end{codebox}

\section{Authentification (\texttt{AuthScreen.js})}

Firebase Authentication gère la connexion et l'inscription. Les messages d'erreur sont traduits en français pour améliorer l'UX. Persistance via \texttt{AsyncStorage}.

\begin{codebox}[listing options={style=jsstyle,language=bash}]
npm install @react-native-async-storage/async-storage
npm install @expo/vector-icons expo-linear-gradient
\end{codebox}

\begin{stepbox}[title={\faLock~Fonctions clés}]
\begin{itemize}
  \item \texttt{handleLogin} \textemdash{} connexion via email/mot de passe
  \item \texttt{handleSignUp} \textemdash{} création de compte et redirection
  \item \texttt{translateError} \textemdash{} mapping Firebase $\rightarrow$ français
  \item \texttt{loadFonts} + \texttt{prepare} \textemdash{} splash screen et polices personnalisées
\end{itemize}
\end{stepbox}

\section{Navigation principale (\texttt{Navigation.js})}

Une navigation \textbf{stack} empile les écrans (\texttt{AuthScreen}, \texttt{SwipeScreen}, \texttt{ChatScreen}, etc.). \texttt{react-native-screens} optimise les transitions natives.

\begin{codebox}[listing options={style=jsstyle,language=bash}]
npm install @react-navigation/native @react-navigation/stack \
  react-native-screens react-native-safe-area-context \
  react-native-gesture-handler react-native-reanimated react-native-svg
\end{codebox}

\section{Système de swipe (\texttt{SwipeScreen.js})}

Cœur de l'app : \texttt{PanResponder} détecte les gestes de glissement et déclenche des animations selon la direction. \texttt{expo-av} joue les extraits Spotify.

\begin{stepbox}[title={\faHandPointer~Logique de geste}]
\begin{itemize}
  \item Swipe droite ($\geq$ seuil) $\rightarrow$ ajout aux \texttt{likedTracks}
  \item Swipe gauche ($\leq$ -seuil) $\rightarrow$ ajout aux \texttt{dislikedTracks}
  \item Geste insuffisant $\rightarrow$ animation de retour à la position initiale
\end{itemize}
\end{stepbox}

\textbf{Dépendances :}
\begin{itemize}
  \item \texttt{react-native-gesture-handler} \textemdash{} gestion des swipes
  \item \texttt{react-native-deck-swiper} \textemdash{} composant deck-like
  \item \texttt{expo-av} \textemdash{} lecture audio des extraits
\end{itemize}

\section{Chat temps réel (\texttt{ChatScreen.js})}

Firestore synchronise les messages en temps réel via \texttt{onSnapshot}. Chaque conversation est identifiée par un \texttt{chatId}, et chaque message inclut l'identité de l'expéditeur et un timestamp.

\begin{codebox}
import { onSnapshot, collection, addDoc, query, orderBy } from "firebase/firestore";

const q = query(
  collection(db, "chats", chatId, "messages"),
  orderBy("timestamp", "asc")
);
const unsubscribe = onSnapshot(q, (snapshot) => {
  setMessages(snapshot.docs.map(d => ({ id: d.id, ...d.data() })));
});
\end{codebox}

\section{Recherche musicale (\texttt{SearchScreen.js})}

Une fonction de normalisation des chaînes (\texttt{normalizeString}) standardise la requête utilisateur (accents, casse). Les résultats sont triés par pertinence et popularité, limités aux 20 premiers, filtrés pour n'inclure que les morceaux avec extrait audio disponible.

\textbf{Dépendances :} \texttt{expo-network} (vérifie la connexion avant requête) et \texttt{buffer} (encodage des données API).

\section{Historique : Liked / Disliked}

Deux écrans miroirs (\texttt{LikedTracksScreen.js}, \texttt{DislikedTracksScreen.js}) affichent les morceaux via \texttt{FlatList}. Chaque ligne montre titre, artiste, album \textemdash{} avec \texttt{"inconnu"} en cas de données manquantes.

\section{Intégration Spotify (\texttt{SpotifyService.js})}

Service centralisé pour toutes les requêtes Spotify. L'authentification se fait via \textbf{client credentials} : \texttt{clientId} + \texttt{clientSecret} configurés dans la console Spotify Developer.

\begin{codebox}
async function getAccessToken() {
  const credentials = base64.encode(`${clientId}:${clientSecret}`);
  const res = await fetch("https://accounts.spotify.com/api/token", {
    method: "POST",
    headers: {
      "Authorization": `Basic ${credentials}`,
      "Content-Type": "application/x-www-form-urlencoded"
    },
    body: "grant_type=client_credentials"
  });
  const data = await res.json();
  return data.access_token;
}
\end{codebox}

Deux fonctions principales :
\begin{itemize}
  \item \texttt{getRandomTracksByGenre(genre)} \textemdash{} alimente le swipe (défaut : \og hip-hop \fg{} + \og french \fg{})
  \item \texttt{searchTracks(query)} \textemdash{} alimente la recherche utilisateur
\end{itemize}

Les deux filtrent pour ne retenir que les morceaux disposant d'\texttt{preview\_url} non null.

\section{Configuration Firebase (\texttt{firebaseConfig.js})}

\texttt{initializeAuth} avec persistance \texttt{AsyncStorage} garantit que la session survit aux redémarrages de l'app. \texttt{getCurrentUser} fournit aux écrans l'identité de l'utilisateur connecté pour personnaliser l'expérience.

\section{Récapitulatif des dépendances}

\renewcommand{\arraystretch}{1.25}
\begin{tabularx}{\linewidth}{@{}>{\bfseries\ttfamily}l X@{}}
\toprule
\rowcolor{light} \normalfont\textbf{Package} & Rôle \\ \midrule
expo                      & SDK React Native \\
firebase                  & Auth + Firestore \\
@react-navigation/*       & Navigation entre écrans \\
react-native-gesture-handler & Gestion des swipes \\
react-native-deck-swiper  & Composant deck \\
expo-av                   & Lecture audio des previews \\
expo-linear-gradient      & Dégradés UI \\
expo-font, expo-splash-screen & Polices, splash screen \\
@react-native-async-storage/async-storage & Persistance locale \\
spotify-web-api-node      & Client Spotify \\
axios                     & Requêtes HTTP \\
react-native-base64       & Encodage clés API \\
\bottomrule
\end{tabularx}

\section{Difficultés rencontrées}

\begin{tabularx}{\linewidth}{@{}>{\bfseries}p{4cm} X@{}}
\toprule
\rowcolor{light} Problème & Solution apportée \\ \midrule
\texttt{react-native-reanimated} exige une config native supplémentaire & Suivi de la doc officielle, configuration Babel \\
Hétérogénéité des machines de l'équipe & Synchronisation via GitHub, sessions partagées \\
Gestion des tokens Spotify (expiration, scopes) & Itérations dans Spotify Developer Console, refresh à la demande \\
\bottomrule
\end{tabularx}

\section{Conclusion}

Harmely combine des briques modernes du développement mobile : framework cross-platform, backend serverless, API tierce, gestures naturels. La SAE a permis de couvrir le cycle complet d'une app : conception UX, intégration API, persistance, temps réel, et déploiement.

\medskip
Au-delà du livrable, l'expérience confirme une réalité du métier : l'essentiel du travail consiste à \textbf{intégrer} des services existants et à soigner l'expérience utilisateur \textemdash{} pas à réécrire l'OAuth ou la lecture audio.

\vspace{0.5cm}
\begin{center}
{\sffamily\itshape\color{midgray}\large
\og La meilleure interface est celle qu'on ne voit pas. \fg{}}
\end{center}

\end{document}
